% Figure: the fast-slow control architecture.
% Included from main.tex. The preamble of main.tex loads the tikz package and
% the arrows.meta, positioning, fit and backgrounds libraries.
\begin{figure}[t]
\centering
\resizebox{\textwidth}{!}{%
\begin{tikzpicture}[
    font=\footnotesize,
    node distance=5mm and 6mm,
    every node/.style={align=center},
    stage/.style={draw, rounded corners=2pt, inner sep=6pt, fill=black!5},
    cluster/.style={draw, rounded corners=3pt},
    flow/.style={-{Stealth[length=2.2mm]}, thick},
    back/.style={-{Stealth[length=2.2mm]}, thick, dashed},
    note/.style={draw, rounded corners=2pt, inner sep=5pt, fill=black!3, font=\scriptsize},
    badge/.style={circle, draw, inner sep=1.2pt, font=\tiny, fill=white},
]
% --- row 1: the fast loop -------------------------------------------
\node[stage] (problem) {Variational problem\\[1pt] $H$, \quad ansatz $U(\vec{\theta})$};
\node[badge, above=0.5mm of problem] {1};
\node[stage, right=of problem] (theta) {$\vec{\theta}_k$};
\node[stage, right=of theta] (spsa) {SPSA step\\[1pt] $\vec{\theta}_{k+1} = \vec{\theta}_k - a_k \hat{g}_k$};
\node[stage, right=of spsa] (energy) {$E(\vec{\theta}_k) = \langle \psi(\vec{\theta}_k) \rvert H \lvert \psi(\vec{\theta}_k) \rangle$};
\begin{scope}[on background layer]
  \node[cluster, fill=black!10, fit=(theta)(spsa)(energy), inner sep=8pt,
        label={[font=\scriptsize\itshape]above:fast loop, $N_w$ iterations at local speed}] (fast) {};
\end{scope}
\node[badge, above=0.5mm of fast] {2};
\draw[flow] (problem) -- (theta);
\draw[flow] (theta) -- (spsa);
\draw[flow] (spsa) -- (energy);
\draw[flow] (energy.north) .. controls +(up:5mm) and +(up:5mm) .. node[above, font=\scriptsize] {$k = 1, \ldots, N_w$} (theta.north);

\node[stage, right=9mm of energy] (window) {Telemetry window $T_k$\\[1pt]
  $\bigl[ E_{k - N_w : k}, \, \lVert \nabla E \rVert_2, \, \eta_k, \, \mathrm{Var}(\vec{\theta}) \bigr]$};
\node[badge, above=0.5mm of window] {3};
\draw[flow] (energy) -- (window);

% --- warm start ------------------------------------------------------
\coordinate (wsa) at ([yshift=8mm]problem.north);
\coordinate (wsb) at ([yshift=8mm]fast.north -| window);
\draw[back] (problem.north) -- (wsa) -| (wsb) node[above, font=\scriptsize, pos=0.4]
  {epoch $0$: proposed $\vec{\theta}^{(0)}$ from the symbolic description of $H$};

% --- row 2: the slow loop -------------------------------------------
\node[cluster, fill=black!10, below=12mm of fast.south -| window, inner sep=8pt] (llm) {%
  LLM meta-optimizer (slow loop)\\[2pt]
  \begin{tabular}{@{}r@{\ }l@{}}
    input: & system prompt + telemetry window \\
    output: & diagnosis, justification, expected effect \\
    format: & one JSON object, validated by a schema \\
  \end{tabular}};
\node[badge, left=1mm of llm] {4};
\draw[flow] (window.south) -- (llm.north -| window);

\node[stage, right=9mm of llm] (action) {Intervention $a$\\[1pt]
  $\eta$ scaled, $\sigma$ noise, restart};
\node[badge, above=0.5mm of action] {5};
\draw[flow] (llm) -- (action);

% --- feedback --------------------------------------------------------
\draw[back] (action.north) -- ++(0,4mm) -| node[above, font=\scriptsize, pos=0.35]
  {applied to $\vec{\theta}$ every $N_w$ iterations} (theta.south);

% --- audit strip -----------------------------------------------------
\node[note, below=7mm of llm.south -| spsa, xshift=-14mm] (latency)
  {latency accounting: $t_{\mathrm{call}}$ measured per call, failed calls included in the cost};
\node[note, right=6mm of latency] (fidelity)
  {XAI audit: declared effect vs.\ observed effect, diagnosis vs.\ ground-truth regime};
\end{tikzpicture}%
}
\caption{The fast-slow control architecture. A stochastic local optimizer
improves the variational angles for $N_w$ iterations between two slow-loop
queries. The controller receives a telemetry window and returns a diagnosis, a
justification, a predicted effect and one intervention. The intervention is
limited to a scaling of the learning rate, a Gaussian perturbation of the
angles and a restart, and it is applied to the angles without modifying the
ansatz.}
\label{fig:methodology}
\end{figure}
